\documentclass[10pt]{article}
\usepackage[margin=0.9in]{geometry}
\usepackage{amsmath,amssymb,booktabs,microtype,xcolor,fancyhdr}
\usepackage[hidelinks]{hyperref}
\definecolor{navy}{RGB}{24,55,95}
\pagestyle{fancy}\fancyhf{}\lhead{$G_2$ on $V_7$}\rhead{Proof and exact verification}\cfoot{\thepage}
\setlength{\headheight}{14pt}
\newcommand{\Sg}{\mathcal S}
\newcommand{\CT}{\operatorname{CT}}
\title{\color{navy}\textbf{The compact $G_2$ sigma-MGF on $V_7$}\\[1mm]\large Proof and exact low-degree matrix verification}
\author{Computational proof companion}\date{17 July 2026}
\begin{document}\maketitle
\begin{abstract}
We prove the compact Molien--Weyl formula for the minimal seven-dimensional
$G_2$-module, construct its maximal-torus matrices from the weights, and
verify every Schur coefficient through degree three exactly.  The result is
$\Sg_{G_2,7}=1+h_2+e_3+O_{\ge4}$, and the ordinary specialization is
$(1-t^2)^{-1}$.
\end{abstract}

\section{General compact-group identity}
For a unitary $K$-module $V$ and eigenvalues $\lambda_j(g)$, the Cauchy
identity gives, as a formal power series,
\[
 \prod_{a,j}(1-t_a\lambda_j(g))^{-1}
 =\sum_\lambda s_\lambda(\mathbf t)\,\chi_{\mathbb S_\lambda V}(g).
\]
Haar integration projects each representation onto its invariant subspace;
hence
\[
 \boxed{\Sg_{K,V}(\mathbf t)=\int_K\prod_a\det(I-t_a\rho(g))^{-1}\,dg
 =\sum_\lambda\dim(\mathbb S_\lambda V)^K s_\lambda(\mathbf t).}
\]
Expanding each determinant separately also gives the requested monomial form
\[
 [m_\tau]\Sg_{K,V}
 =\dim\left(\bigotimes_i\operatorname{Sym}^{\tau_i}V\right)^K.
\]
If $T$ is a maximal torus, Weyl integration yields
\[
 \Sg_{K,V}=\frac1{|W|}\CT_z\!\left[
 \prod_{\alpha>0}(1-z^\alpha)(1-z^{-\alpha})
 \prod_{a,\mu}(1-t_a z^\mu)^{-m(\mu)}\right].
\]
This proves the proposed compact analogue of the finite Molien formula.

\section{Specialization to the minimal module}
The weights of $V_7$ are zero together with the six short roots.  Since
$|W(G_2)|=12$,
\[
 \boxed{\Sg_{G_2,7}=\frac1{12}\CT_z\left[D_{G_2}(z)
 \prod_a\left((1-t_a)^{-1}\prod_{\beta\in\Phi_{\rm short}}
 (1-t_a z^\beta)^{-1}\right)\right].}
\]
The invariant metric lies in $\operatorname{Sym}^2V_7$, while the octonionic
associative form lies in $\bigwedge^3V_7$.  Their uniqueness and the absence
of other invariants below degree four give
\[
 \boxed{\Sg_{G_2,7}=1+h_2+e_3+O_{\ge4}.}
\]
A standard invariant-theory theorem for the minimal $G_2$ module states that
$\mathbb C[V_7]^{G_2}=\mathbb C[q]$, where $q$ is the invariant quadratic
norm.  Using that external theorem,
$\Sg_{G_2,7}(t,0,\ldots)=(1-t^2)^{-1}$.  The exact Weyl computation in this
note independently verifies the coefficients through degree three.

\section{Exact matrix verification}
The notebook generates the root system from the $G_2$ Cartan matrix and forms
\[
 M(\theta)=\operatorname{diag}\bigl(e^{i\langle\mu,\theta\rangle}\bigr)_{
 \mu\in\{0\}\cup\Phi_{\rm short}}\in U(7).
\]
It compares $\prod_a\det(I-t_aM)^{-1}$ directly with the product over the
seven diagonal eigenvalues.  For integration it evaluates Schur characters by
Jacobi--Trudi and uses the one-sided Weyl denominator
\[
 \prod_{\alpha>0}(1-z^{-\alpha})
 =\sum_{w\in W}\det(w)z^{w\rho-\rho}.
\]
Thus every coefficient is an exact integer, with no quadrature tolerance.
In the requested monomial basis the exact degree-at-most-three signature is
\[
([m_{(2)}],[m_{(1,1)}],[m_{(3)}],[m_{(2,1)}],[m_{(1,1,1)}])
=(1,1,0,0,1).
\]

\begin{center}\small
\begin{tabular}{@{}c|rrrrrrr@{}}\toprule
$\lambda$&$\varnothing$&$(1)$&$(2)$&$(1,1)$&$(3)$&$(2,1)$&$(1,1,1)$\\\midrule
proposed&1&0&1&0&0&0&1\\
exact Weyl check&1&0&1&0&0&0&1\\\bottomrule
\end{tabular}
\end{center}
The sampled torus matrix has unitarity residual $3.2\times10^{-16}$ and the
two integrand evaluations differ by less than $9\times10^{-16}$.  All seven
coefficient checks pass.

\paragraph{Reproducibility.}
Run \texttt{marimo\_notebooks/lie\_group\_g2.py} with marimo.  The shared exact-arithmetic
module is in the parent folder.

\section*{External invariant-theory source}
See J.~F.~Adams, \emph{Lectures on Exceptional Lie Groups}, Chicago
Lectures in Mathematics (1996), for the minimal $G_2$ module and its invariant
metric.  The all-degree one-copy invariant-ring statement is an external
input; the low-degree Weyl checks above are internal to this artifact.
\end{document}
